\documentclass[12pt]{article}
\usepackage[utf8]{inputenc}
\usepackage[T1]{fontenc}
\usepackage{amsmath,amssymb,amscd}
\usepackage{geometry}
\geometry{margin=2.5cm}

\begin{document}

\begin{center}
\fbox{\textbf{CURS 6}}\\[2pt]
\underline{(ALGEBRA)}
\end{center}
\begin{flushright}
\textit{8-XI-94}
\end{flushright}

\bigskip
\noindent\underline{\textbf{Ordinul unui element}}

\bigskip
\noindent\underline{Anexa I}: -- Elemente de aritmetică în $\mathbb{N}$.

\bigskip
\noindent\underline{Sistem Peano}: $(A,a,s)$ ; $a\in A$; $s:A\to A$,\footnote{în original apare ,,$s:A\to B$''; corectat aici la $A\to A$, cum cere definiția unui sistem Peano.} $A\neq\varnothing$

\begin{enumerate}
\item[1)] $s(x)\neq a$ $(\forall)\,x\in A$, \ $\big(s(A)=A\setminus\{a\}\big)$.
\item[2)] $s$ injectivă.
\item[3)] Dacă $M\subseteq A$ are pp.\ $\begin{cases} a\in M \\ x\in M \Rightarrow s(x)\in M \end{cases} \Rightarrow M=A$.
\end{enumerate}

\bigskip
\noindent\fbox{T1}. \ Dacă $(A,a,s)$, $(B,b,t)$ sisteme Peano $\Rightarrow$

\noindent $\exists!\, f:A\to B$ a.î.\ $f(a)=b$, $f\circ s = t\circ f$.

\noindent Mai mult,\footnote{fraza exactă e neclară în original.} $f$ e bijectivă.

\[
\begin{CD}
A @>s>> A\\
@VfVV @VVfV\\
B @>t>> B
\end{CD}
\]

\bigskip
\noindent Fie $(\mathbb{N},0,s)$; \ $s:\mathbb{N}\to\mathbb{N}$.

\noindent funcţ.\ succesor.

\noindent $s(x)=x'$ -- succesorul lui $x$.

\bigskip
\noindent N.\ Becheanu -- Alg.\ ptr.\ perf.\ profesorale\footnote{referință bibliografică; titlul exact e reconstituit din formă compactată.}

\bigskip
\noindent\fbox{T2}: $\exists!\, e:\mathbb{N}\times\mathbb{N}\to\mathbb{N}$ \ $e\overset{\text{not}}{=}+$

\noindent $A_1$: $e(m,0)=m$ $(\forall)\,m\in\mathbb{N}$.

\noindent $A_2$: $e(m,n')=e(m,n)'$ $(\forall)\,m,n\in\mathbb{N}$ \quad $m+n'=(m+n)'$.

\bigskip
\noindent Rezultă $m+p=m+n \Rightarrow p=n$.

\[
(\mathbb{N},+,0) = \text{monoid comutativ cu simplificare.}
\]

\bigskip
\noindent\fbox{T}: $\exists!\, \psi:\mathbb{N}\times\mathbb{N}\to\mathbb{N}$ a.î.

\noindent $\psi(m,0)=0$\footnote{în original pare ,,$=m$''; corectat la $0$, cum cere definiția standard a înmulțirii (notația alăturată ,,$m\cdot 0$'' e la fel $0$, nu $m$).} $(\forall)\,m\in\mathbb{N}$: \quad $m\cdot 0=0$

\noindent $\psi(m,n')=\psi(m,n)+m$ $(\forall)\,m,n\in\mathbb{N}$ \quad $mn'=mn+m$.

\bigskip
\noindent $(\mathbb{N},\cdot,1)$ = monoid comutativ cu simplificare.

\noindent $mn=mp \Rightarrow n=p$.

\noindent $m(n+p)=mn+mp$.

\bigskip
\noindent $(\mathbb{N},+,\cdot)$ = structură algebrică.

\noindent $(\mathbb{N},\leq)$ = structură de ordine.

\end{document}
